\section{三角置換}\label{節:三角置換の初歩}

前節では根号そのものを $t$ と置いた.しかし $\sqrt{a^2-x^2}$ $(a>0)$ では,\,$t=\sqrt{a^2-x^2}$ と置いて $x$ を解き直しても,\,$x=\pm\sqrt{a^2-t^2}$ となって根号が残ってしまう.では,根号の中の平方の差を消すにはどうすればよいだろうか.

ここで三角関数の関係式
\[
 1-\sin^2\theta=\cos^2\theta,\qquad
 1+\tan^2\theta=\frac1{\cos^2\theta}
\]
を思い出そう.\,$a^2-x^2$ は前者と,\,$a^2+x^2$ は後者と同じ形をしている.そこで $x$ を三角関数で表し,根号や分母を $\cos\theta$ の式に変えて,前章の三角関数の積分に帰着させることを目標とする.

まず $\sqrt{a^2-x^2}$ である.\,$x=a\sin\theta$ と置くと
\begin{align*}
 \sqrt{a^2-x^2}
 &=a\sqrt{1-\sin^2\theta}\\
 &=a\sqrt{\cos^2\theta}\\
 &=a|\cos\theta|
\end{align*}
となる.最後は $a\cos\theta$ ではなく,まず絶対値付きである.\,$\sqrt{z^2}=|z|$ だからである.そこで $\theta$ の範囲として $-\frac\pi2\le\theta\le\frac\pi2$ を選ぶ.この範囲では $\cos\theta\ge0$ となり,絶対値を外せる.さらにこの範囲では,図\ref{図:正弦の枝}のように $\sin\theta$ が狭義単調増加なので,\,$-a\le x\le a$ の各 $x$ に対して $\theta$ が1つに決まる.前節で述べた枝を選んだことになる.このとき
\[
 \sqrt{a^2-x^2}=a\cos\theta,\qquad
 \dd x=a\cos\theta\dd\theta
\]
であり,\,$\sqrt{a^2-x^2}\dd x=a^2\cos^2\theta\dd\theta$ となる.あとは二倍角の公式で積分できる.
次に $a^2+x^2$ である.\,$x=a\tan\theta$ と置き,\,$-\frac\pi2<\theta<\frac\pi2$ を選ぶ.この範囲で正接は狭義単調増加し,すべての実数 $x$ に $\theta$ が1つずつ対応する.このとき
\[
 a^2+x^2=a^2(1+\tan^2\theta)=\frac{a^2}{\cos^2\theta},\qquad
 \dd x=\frac a{\cos^2\theta}\dd\theta
\]
となる.分母と微分に同じ $\frac1{\cos^2\theta}$ が現れるので,
\begin{align*}
 \frac1{a^2+x^2}\dd x
 &=\frac{\cos^2\theta}{a^2}\cdot\frac a{\cos^2\theta}\dd\theta\\
 &=\frac1a\dd\theta
\end{align*}
と,被積分関数が定数になる.

\begin{figure}[htbp]
 \centering
 \begin{tikzpicture}[x=2cm,y=1.5cm]
  \draw[-{Stealth}] (-1.95,0)--(1.95,0) node[right] {$\theta$};
  \draw[-{Stealth}] (0,-1.3)--(0,1.35) node[above] {$\sin\theta$};
  \draw[thick,blue,domain=-1.570796:1.570796,samples=70]
    plot (\x,{sin(\x r)});
  \draw[dashed] (-1.570796,0)--(-1.570796,-1)--(0,-1);
  \draw[dashed] (1.570796,0)--(1.570796,1)--(0,1);
  \node[above] at (-1.570796,0) {$-\frac\pi2$};
  \node[below] at (1.570796,0) {$\frac\pi2$};
  \node[above left] at (0,0) {$0$};
  \node[left] at (0,1) {$1$};
  \node[right] at (0,-1) {$-1$};
 \end{tikzpicture}
 \caption{正弦はこの区間で狭義単調増加し,逆向きの対応を1つに決められる.}
 \label{図:正弦の枝}
\end{figure}

\begin{bookpoint}
\textbf{POINT}\par
$\sqrt{a^2-x^2}$ には $x=a\sin\theta$ $\left(-\frac\pi2\le\theta\le\frac\pi2\right)$,\,$a^2+x^2$ には $x=a\tan\theta$ $\left(-\frac\pi2<\theta<\frac\pi2\right)$ と置き,端点を $\theta$ に移して $\theta$ で積分する.
\end{bookpoint}

\begin{example}{四分円の面積}
$a>0$ とする.\,$\int_0^a\sqrt{a^2-x^2}\dd x$ を求めよ.
\end{example}

\paragraphbf{方針}
$x=a\sin\theta$ と置き,\,$0\le\theta\le\frac\pi2$ を選んで $\cos^2\theta$ の積分に帰着させる.

\paragraphbf{解答}
$\frac{\pi a^2}{4}$

\paragraphbf{解法}
ここでは $x\ge0$ だから,\,$x=a\sin\theta$ と置き,\,$0\le\theta\le\frac\pi2$ を選ぶ.端点の対応は
\[
 \begin{array}{c|ccc}
 x&0&\to&a\\ \hline
 \theta&0&\to&\frac\pi2
 \end{array}
\]
である.この範囲では $\cos\theta\ge0$ なので
\[
 \dd x=a\cos\theta\dd\theta,\qquad
 \sqrt{a^2-x^2}=a\cos\theta
\]
となり,
\begin{align*}
 \int_0^a\sqrt{a^2-x^2}\dd x
 &=a^2\int_0^{\frac\pi2}\cos^2\theta\dd\theta\\
 &=a^2\left[\frac\theta2+\frac14\sin(2\theta)\right]_0^{\frac\pi2}\\
 &=\frac{\pi a^2}{4}
\end{align*}
を得る.曲線 $y=\sqrt{a^2-x^2}$ は円 $x^2+y^2=a^2$ の上半分であり,この積分は半径 $a$ の円の4分の1の面積だから,幾何学的な結果とも一致する.積分計算そのものは,三角置換と二倍角の公式で得られている.

\begin{example}{正接への置換}
$\int_0^1\frac1{1+x^2}\dd x$ を求めよ.
\end{example}

\paragraphbf{方針}
$x=\tan\theta$ と置き,\,$0\le\theta\le\frac\pi4$ を選んで,被積分関数を定数にする.

\paragraphbf{解答}
$\frac\pi4$

\paragraphbf{解法}
$x=\tan\theta$ と置き,\,$0\le\theta\le\frac\pi4$ を選ぶ.この区間で正接は狭義単調増加し,端点 $\theta=0,\frac\pi4$ は $x=0,1$ に対応する.\,$\dd x=\frac1{\cos^2\theta}\dd\theta$ だから,
\begin{align*}
 \int_0^1\frac1{1+x^2}\dd x
 &=\int_0^{\frac\pi4}
   \frac1{1+\tan^2\theta}\cdot
   \frac1{\cos^2\theta}\dd\theta\\
 &=\int_0^{\frac\pi4}1\dd\theta\\
 &=\frac\pi4
\end{align*}
を得る.

\paragraphbf{$\theta$ の範囲を選ぶ理由}
四分円の例題で $0\le\theta\le\pi$ を選ぶと,\,$\sin\theta$ は $0$ から $1$ へ増えた後,\,$0$ へ戻る.この区間全体を $0\le x\le a$ と1対1に対応させることはできず,後半では $\cos\theta$ も負になる.枝を指定することは,こうした重複と符号を確定するために必要なのである.

\paragraphbf{不定積分にする場合}
正接の例題を不定積分として表すなら,最後に $\tan\theta=x$ を満たす $\theta$ を $x$ の関数として書くことになる.この逆関数の名前,枝,微分は\chapterref{章:置換の背景}で詳しく扱う.ここでは定積分の端点が既知の角に対応するため,逆関数の記号を導入しなくても計算を終えられる.

\subsection{図形と置換}
定積分$\int_0^1\sqrt{1-x^2}\dd x$は,半径$1$の四分円の面積だから$\frac\pi4$である.三角置換でも確かめてみよう.\,$x=\sin\theta$,\,$0\le\theta\le\frac\pi2$とすると,
\[
 \int_0^1\sqrt{1-x^2}\dd x
 =\int_0^{\frac\pi2}\cos^2\theta\dd\theta
 =\frac12\int_0^{\frac\pi2}(1+\cos2\theta)\dd\theta=\frac\pi4.
\]
この場合は図形から値を求めるほうが短い.しかし,同じ置換は$\int_0^a\sqrt{1-x^2}\dd x$にも使え,途中までの面積を表せる.変数$x$を円上の角$\theta$で表したのである.その際には,高さの$\cos\theta$に加え,横幅の変化からも$\dd x=\cos\theta\dd\theta$が現れる.円上の高さを置き換えるだけでは,面積の計算にはならない.

%
